\section{有理函数的不定积分、可化为有理函数的不定积分}\label{sec:5.3}

由两个实多项式函数之商所表示的函数称为 (实) 有理函数, 它的一般形式为
\[
R(x)=\frac{P(x)}{Q(x)}
=\frac{\alpha_{0}x^{n}+\alpha_{1}x^{n-1}+\cdots+\alpha_{n-1}x+\alpha_{n}}
{\beta_{0}x^{m}+\beta_{1}x^{m-1}+\cdots+\beta_{m-1}x+\beta_{m}},
\]
其中 $n,m$ 为非负整数, $\alpha_{0},\alpha_{1},\cdots,\alpha_{n}$ 与 $\beta_{0},\beta_{1},\cdots,\beta_{m}$ 都为实常数, 且 $\alpha_{0}\neq 0$, $\beta_{0}\neq 0$.

如果 $m>n$, 则称它为真分式; 如果 $m\leqslant n$, 则称它为假分式. 由多项式除法可知,
\[
R(x)=\frac{P(x)}{Q(x)}=\ \text{多项式}\ +\ \text{真分式}.
\]
根据不定积分的简单性质及 $\frac{x^{k+1}}{k+1}$ 为 $x^{k}$ 的原函数可知, 多项式函数的不定积分是容易得到的. 因此, 要计算有理函数的不定积分只需研究真分式的不定积分.

在代数学中, 有部分分式分解定理:

设
\[
R(x)=\frac{P(x)}{Q(x)}
=\frac{\alpha_{0}x^{n}+\alpha_{1}x^{n-1}+\cdots+\alpha_{n-1}x+\alpha_{n}}
{\beta_{0}x^{m}+\beta_{1}x^{m-1}+\cdots+\beta_{m-1}x+\beta_{m}}
\]
为真分式, $\alpha_{0}\neq 0$, $\beta_{0}\neq 0$, $n<m$, 其中
\[
Q(x)=\beta_{0}(x-a_{1})^{\lambda_{1}}\cdots(x-a_{s})^{\lambda_{s}}
(x^{2}+p_{1}x+q_{1})^{\mu_{1}}\cdots(x^{2}+p_{t}x+q_{t})^{\mu_{t}},
\]
这里 $\lambda_{i},\mu_{j}$ ($i=1,2,\cdots,s$; $j=1,2,\cdots,t$) 均为自然数, 且
\[
\sum_{i=1}^{s}\lambda_{i}+2\sum_{j=1}^{t}\mu_{j}=m;\quad p_{j}^{2}-4q_{j}<0,\ j=1,2,\cdots,t.
\]
则 $R(x)$ 可惟一分解为若干部分分式 (称为最简分式) 之和 (即部分分式分解):
\begin{align*}
R(x)&=\frac{P(x)}{Q(x)}
=\frac{A_{1}}{x-a_{1}}+\cdots+\frac{A_{\lambda_{1}}}{(x-a_{1})^{\lambda_{1}}}+\cdots+\frac{B_{1}}{x-a_{s}}\\
&\quad+\cdots+\frac{B_{\lambda_{s}}}{(x-a_{s})^{\lambda_{s}}}
+\frac{K_{1}x+L_{1}}{x^{2}+p_{1}x+q_{1}}+\cdots+\frac{K_{\mu_{1}}x+L_{\mu_{1}}}{(x^{2}+p_{1}x+q_{1})^{\mu_{1}}}\\
&\quad+\cdots+\frac{M_{1}x+N_{1}}{x^{2}+p_{t}x+q_{t}}+\cdots+\frac{M_{\mu_{t}}x+N_{\mu_{t}}}{(x^{2}+p_{t}x+q_{t})^{\mu_{t}}},
\end{align*}
其中 $A_{i},B_{i},K_{i},L_{i},M_{i},N_{i}$ 均为实数.

应用这个定理, 对任何真分式 $R(x)$ 的不定积分, 一旦完成了 $R(x)$ 的部分分式的分解, 最后转化为求各个部分分式 (最简分式) 的不定积分. 由上述讨论知, 任何有理真分式的不定积分都将归为求以下两种形式的不定积分:

(1)
\[
\int\frac{\mathrm{d}x}{(x-a)^{k}}=
\begin{cases}
\ln|x-a|+C, & k=1,\\[2mm]
\dfrac{1}{(1-k)(x-a)^{k-1}}+C, & k>1.
\end{cases}
\]

(2)
\begin{align*}
&\int\frac{Kx+L}{(x^{2}+px+q)^{k}}\mathrm{d}x\quad (p^{2}-4q<0)\\
&\quad=\int\frac{K\left(x+\frac{p}{2}\right)+\left(L-\frac{Kp}{2}\right)}
{\left[\left(x+\frac{p}{2}\right)^{2}+\left(q-\frac{p^{2}}{4}\right)\right]^{k}}\mathrm{d}\left(x+\frac{p}{2}\right)\\
&\quad\xlongequal{t=x+\frac{p}{2}}\int\frac{Kt+\left(L-\frac{Kp}{2}\right)}{(t^{2}+r^{2})^{k}}\mathrm{d}t\\
&\quad=K\int\frac{t}{(t^{2}+r^{2})^{k}}\mathrm{d}t+\left(L-\frac{Kp}{2}\right)\int\frac{\mathrm{d}t}{(t^{2}+r^{2})^{k}}.
\end{align*}

当 $k=1$ 时, 上式右边两个不定积分分别为
\[
\int\frac{t}{t^{2}+r^{2}}\mathrm{d}t=\frac{1}{2}\int\frac{\mathrm{d}(t^{2}+r^{2})}{t^{2}+r^{2}}=\frac{1}{2}\ln(t^{2}+r^{2})+C,
\]
\[
\int\frac{\mathrm{d}t}{t^{2}+r^{2}}=\frac{1}{r}\arctan\frac{t}{r}+C.
\]
当 $k\geqslant 2$ 时, 右边第一个积分为
\begin{align*}
\int\frac{t}{(t^{2}+r^{2})^{k}}\mathrm{d}t&=\frac{1}{2}\int\frac{\mathrm{d}(t^{2}+r^{2})}{(t^{2}+r^{2})^{k}}\\
&=\frac{1}{2(1-k)(t^{2}+r^{2})^{k-1}}+C.
\end{align*}
对于第二个不定积分, 记
\[
I_{k}=\int\frac{\mathrm{d}t}{(t^{2}+r^{2})^{k}},
\]
则
\begin{align*}
I_{k}&=\int\frac{\mathrm{d}t}{(t^{2}+r^{2})^{k}}=\frac{1}{r^{2}}\int\frac{(t^{2}+r^{2})-t^{2}}{(t^{2}+r^{2})^{k}}\mathrm{d}t\\
&=\frac{1}{r^{2}}I_{k-1}-\frac{1}{r^{2}}\int\frac{t^{2}}{(t^{2}+r^{2})^{k}}\mathrm{d}t\\
&=\frac{1}{r^{2}}I_{k-1}+\frac{1}{2r^{2}(k-1)}\int t\,\mathrm{d}\frac{1}{(t^{2}+r^{2})^{k-1}}\\
&\xlongequal{\text{分部积分}}\frac{1}{r^{2}}I_{k-1}+\frac{1}{2r^{2}(k-1)}
\left[\frac{t}{(t^{2}+r^{2})^{k-1}}-I_{k-1}\right],
\end{align*}
\[
I_{k}=\frac{t}{2r^{2}(k-1)(t^{2}+r^{2})^{k-1}}+\frac{2k-3}{2r^{2}(k-1)}I_{k-1}\quad(\text{递推公式}).
\]
或者
\begin{align*}
I_{k-1}&=\int\frac{\mathrm{d}t}{(t^{2}+r^{2})^{k-1}}\\
&\xlongequal{\text{分部积分}}\frac{t}{(t^{2}+r^{2})^{k-1}}-\int t\left[-(k-1)\right]\frac{2t}{(t^{2}+r^{2})^{k}}\mathrm{d}t\\
&=\frac{t}{(t^{2}+r^{2})^{k-1}}+2(k-1)\int\frac{(t^{2}+r^{2})-r^{2}}{(t^{2}+r^{2})^{k}}\mathrm{d}t\\
&=\frac{t}{(t^{2}+r^{2})^{k-1}}+2(k-1)I_{k-1}-2r^{2}(k-1)I_{k},
\end{align*}
\[
I_{k}=\frac{t}{2r^{2}(k-1)(t^{2}+r^{2})^{k-1}}+\frac{2k-3}{2r^{2}(k-1)}I_{k-1}.
\]
根据递推公式, 从 $I_{1}=\int\frac{\mathrm{d}t}{t^{2}+r^{2}}=\frac{1}{r}\arctan\frac{t}{r}+C$ 得到 $I_{2}$, 再从 $I_{2}$ 得到 $I_{3}$, $\cdots$, 从 $I_{k-1}$ 得到 $I_{k}$.

注意, 最后需将 $t=x+\frac{p}{2}$ 代回去.

\begin{example}\label{ex:5.3.1}
求 $\int\frac{x^{2}+1}{(x^{2}-2x+2)^{2}}\mathrm{d}x$.
\end{example}

\begin{solution}
\begin{align*}
\int\frac{x^{2}+1}{(x^{2}-2x+2)^{2}}\mathrm{d}x
&=\int\frac{(x^{2}-2x+2)+(2x-1)}{(x^{2}-2x+2)^{2}}\mathrm{d}x\\
&=\int\frac{\mathrm{d}x}{x^{2}-2x+2}+\int\frac{2x-2}{(x^{2}-2x+2)^{2}}\mathrm{d}x+\int\frac{\mathrm{d}x}{(x^{2}-2x+2)^{2}}\\
&=\int\frac{\mathrm{d}(x-1)}{(x-1)^{2}+1}+\int\frac{\mathrm{d}(x^{2}-2x+2)}{(x^{2}-2x+2)^{2}}+\int\frac{\mathrm{d}(x-1)}{\left[(x-1)^{2}+1\right]^{2}}\\
&=\arctan(x-1)-\frac{1}{x^{2}-2x+2}+\frac{x-1}{2(x^{2}-2x+2)}\\
&\quad+\frac{1}{2}\arctan(x-1)+C\\
&=\frac{3}{2}\arctan(x-1)+\frac{x-3}{2(x^{2}-2x+2)}+C,
\end{align*}
其中
\begin{align*}
I_{2}&=\int\frac{\mathrm{d}t}{(t^{2}+1)^{2}}
\xlongequal[k=2,\ r=1]{\text{递推公式}}\frac{t}{2(t^{2}+1)}+\frac{1}{2}\int\frac{\mathrm{d}t}{t^{2}+1}\\
&=\frac{x-1}{2\left[(x-1)^{2}+1\right]}+\frac{1}{2}\arctan(x-1)+C_{1}\\
&=\frac{x-1}{2(x^{2}-2x+2)}+\frac{1}{2}\arctan(x-1)+C_{1}.
\end{align*}
\end{solution}

\begin{example}\label{ex:5.3.2}
求 $\int\frac{5x+6}{x^{2}+x+1}\mathrm{d}x$.
\end{example}

\begin{solution}
\begin{align*}
\int\frac{5x+6}{x^{2}+x+1}\mathrm{d}x&=\frac{5}{2}\int\frac{2x+\frac{12}{5}}{x^{2}+x+1}\mathrm{d}x
=\frac{5}{2}\int\frac{(2x+1)+\frac{7}{5}}{x^{2}+x+1}\mathrm{d}x\\
&=\frac{5}{2}\int\frac{\mathrm{d}(x^{2}+x+1)}{x^{2}+x+1}
+\frac{7}{2}\int\frac{\mathrm{d}\left(x+\frac{1}{2}\right)}{\left(x+\frac{1}{2}\right)^{2}+\left(\frac{\sqrt{3}}{2}\right)^{2}}\\
&=\frac{5}{2}\ln(x^{2}+x+1)+\frac{7}{2}\cdot\frac{1}{\frac{\sqrt{3}}{2}}\arctan\frac{x+\frac{1}{2}}{\frac{\sqrt{3}}{2}}+C\\
&=\frac{5}{2}\ln(x^{2}+x+1)+\frac{7}{\sqrt{3}}\arctan\frac{2x+1}{\sqrt{3}}+C.
\end{align*}
\end{solution}

\begin{example}\label{ex:5.3.3}
求 $\int\frac{x+1}{x^{2}-4x+3}\mathrm{d}x$.
\end{example}

\begin{solution}[解法 1]
\begin{align*}
\int\frac{x+1}{x^{2}-4x+3}\mathrm{d}x&=\int\frac{x+1}{(x-3)(x-1)}\mathrm{d}x\\
&\xlongequal{\text{凑}}\int\left(\frac{2}{x-3}+\frac{-1}{x-1}\right)\mathrm{d}x\\
&=2\ln|x-3|-\ln|x-1|+C=\ln\left|\frac{(x-3)^{2}}{x-1}\right|+C.
\end{align*}
\end{solution}

\begin{solution}[解法 2 (待定系数法)]
设
\[
\frac{x+1}{x^{2}-4x+3}=\frac{A}{x-3}+\frac{B}{x-1},
\]
则
\[
x+1=A(x-1)+B(x-3)=(A+B)x-(A+3B),\tag{$*$}
\]
比较上面等式两边关于 $x$ 的同幂次的系数知,
\[
\begin{cases}
A+B=1,\\
A+3B=-1.
\end{cases}
\]
下式减上式得 $2B=-2$, $B=-1$, 所以 $A=1-B=2$.

或者, 令 $x=1$ 代入 $(*)$ 式得 $2=-2B$, $B=-1$. 令 $x=3$ 代入 $(*)$ 式得 $4=2A$, $A=2$. 代入即可得到与解法 1 相同的结果.
\end{solution}

\begin{example}\label{ex:5.3.4}
求 $\int\frac{5x+3}{(x^{2}-2x+5)^{2}}\mathrm{d}x$.
\end{example}

\begin{solution}[解法 1]
\begin{align*}
\int&\frac{5x+3}{(x^{2}-2x+5)^{2}}\mathrm{d}x\\
&=\int\frac{\frac{5}{2}(2x-2)}{(x^{2}-2x+5)^{2}}\mathrm{d}x+8\int\frac{\mathrm{d}x}{(x^{2}-2x+5)^{2}}\\
&=\frac{5}{2}\int\frac{\mathrm{d}(x^{2}-2x+5)}{(x^{2}-2x+5)^{2}}+8\int\frac{\mathrm{d}(x-1)}{\left[(x-1)^{2}+2^{2}\right]^{2}}\\
&=-\frac{5}{2}\cdot\frac{1}{x^{2}-2x+5}
+8\left[\frac{1}{8}\cdot\frac{x-1}{(x-1)^{2}+4}+\frac{1}{16}\arctan\frac{x-1}{2}\right]+C\\
&=-\frac{2x-7}{2(x^{2}-2x+5)}+\frac{1}{2}\arctan\frac{x-1}{2}+C.
\end{align*}
\end{solution}

\begin{solution}[解法 2]
\begin{align*}
\int\frac{5x+3}{(x^{2}-2x+5)^{2}}\mathrm{d}x
&\xlongequal{u=x-1}\int\frac{5u+8}{(u^{2}+4)^{2}}\mathrm{d}u\\
&=\frac{5}{2}\int\frac{\mathrm{d}(u^{2}+4)}{(u^{2}+4)^{2}}+2\int\frac{(u^{2}+4)-u^{2}}{(u^{2}+4)^{2}}\mathrm{d}u\\
&=-\frac{5}{2}\cdot\frac{1}{u^{2}+4}+2\int\frac{\mathrm{d}u}{u^{2}+4}+\int u\,\mathrm{d}\frac{1}{u^{2}+4}\\
&=-\frac{5}{2(u^{2}+4)}+2\int\frac{\mathrm{d}u}{u^{2}+4}+\frac{u}{u^{2}+4}-\int\frac{\mathrm{d}u}{u^{2}+4}\\
&=\frac{2u-5}{2(u^{2}+4)}+\frac{1}{2}\arctan\frac{u}{2}+C\\
&=\frac{2x-7}{2(x^{2}-2x+5)}+\frac{1}{2}\arctan\frac{x-1}{2}+C.
\end{align*}
\end{solution}

\begin{example}\label{ex:5.3.5}
求 $\int\frac{x^{5}+x^{3}}{x^{3}+1}\mathrm{d}x$.
\end{example}

\begin{solution}
\begin{align*}
\int\frac{x^{5}+x^{3}}{x^{3}+1}\mathrm{d}x
&=\int\left[(x^{2}+1)-\frac{x^{2}+1}{x^{3}+1}\right]\mathrm{d}x\\
&=\int\left[(x^{2}+1)-\frac{2}{3}\cdot\frac{1}{x+1}-\frac{1}{3}\cdot\frac{x+1}{x^{2}-x+1}\right]\mathrm{d}x\\
&=\int\left(x^{2}+1-\frac{2}{3}\cdot\frac{1}{x+1}-\frac{1}{6}\cdot\frac{2x-1}{x^{2}-x+1}-\frac{1}{2}\cdot\frac{1}{x^{2}-x+1}\right)\mathrm{d}x\\
&=\frac{x^{3}}{3}+x-\frac{1}{6}\ln(x^{3}+1)(x+1)^{3}-\frac{1}{\sqrt{3}}\arctan\frac{2(x-1)}{\sqrt{3}}+C.
\end{align*}
其中 $\frac{x^{5}+x^{3}}{x^{3}+1}=(x^{2}+1)-\frac{x^{2}+1}{x^{3}+1}$ 可用凑的方法, 也可用多项式的辗转相除得到. 而应用待定系数法可令
\[
\frac{x^{2}+1}{x^{3}+1}=\frac{A}{x+1}+\frac{Bx+C}{x^{2}-x+1},
\]
则有
\begin{align*}
x^{2}+1&=A(x^{2}-x+1)+(Bx+C)(x+1)\\
&=(A+B)x^{2}+(-A+B+C)x+(A+C),
\end{align*}
从而有
\[
\left\{\begin{array}
{l@{\qquad}r}
A+B=1, & \text{(1)}\\
-A+B+C=0, & \text{(2)}\\
A+C=1. & \text{(3)}
\end{array}\right.
\]
于是, 由 (1)$+$(2) 及 (2)$+$(3) 得到
\[
\left\{\begin{array}
{l@{\qquad}r}
2B+C=1, & \text{(4)}\\
B+2C=1. & \text{(5)}
\end{array}\right.
\]
再由 $2\times$(5)$-$(4) 得 $3C=1$, $C=\frac{1}{3}$; $B=1-2C=1-\frac{2}{3}=\frac{1}{3}$; $A=1-B=1-\frac{1}{3}=\frac{2}{3}$.
\end{solution}

\begin{example}\label{ex:5.3.6}
求 $\int\frac{x^{3}+1}{x^{4}-3x^{3}+3x^{2}-x}\mathrm{d}x$.
\end{example}

\begin{solution}[解法 1]
$x^{4}-3x^{3}+3x^{2}-x=x(x-1)^{3}$. 用待定系数法, 令
\begin{align*}
\frac{x^{3}+1}{x(x-1)^{3}}&=\frac{x^{3}+1}{x^{4}-3x^{3}+3x^{2}-x}\\
&=\frac{A}{x}+\frac{B}{(x-1)^{3}}+\frac{C}{(x-1)^{2}}+\frac{D}{x-1},
\end{align*}
则
\begin{align*}
x^{3}+1&=A(x-1)^{3}+Bx+Cx(x-1)+Dx(x-1)^{2}\\
&=(A+D)x^{3}+(-3A+C-2D)x^{2}+(3A+B-C+D)x-A,
\end{align*}
于是
\[
\begin{cases}
A+D=1,\\
-3A+C-2D=0,\\
3A+B-C+D=0,\\
-A=1.
\end{cases}\quad
\text{即}\quad
\begin{cases}
A=-1,\\
D=1-A=1-(-1)=2,\\
C=3A+2D=-3+4=1,\\
B=-3A+C-D=3+1-2=2.
\end{cases}
\]
所以
\begin{align*}
&\int\frac{x^{3}+1}{x^{4}-3x^{3}+3x^{2}-x}\mathrm{d}x\\
&=\int\left[\frac{-1}{x}+\frac{2}{(x-1)^{3}}+\frac{1}{(x-1)^{2}}+\frac{2}{x-1}\right]\mathrm{d}x\\
&=-1\cdot\ln|x|-\frac{1}{(x-1)^{2}}-\frac{1}{x-1}+2\ln|x-1|+C\\
&=\ln\frac{(x-1)^{2}}{|x|}-\frac{x}{(x-1)^{2}}+C.
\end{align*}
\end{solution}

\begin{solution}[解法 2]
设 $x^{3}+1=A(x-1)^{3}+Bx+Cx(x-1)+Dx(x-1)^{2}$. 令 $x=1$ 代入上式得 $B=2$; $x=0$ 代入上式得 $A=-1$. 从而
\[
x^{3}+1=-(x-1)^{3}+2x+Cx(x-1)+Dx(x-1)^{2},
\]
\begin{align*}
Cx(x-1)+Dx(x-1)^{2}&=x^{3}+1+x^{3}-3x^{2}+3x-1-2x\\
&=2x^{3}-3x^{2}+x=x(2x^{2}-3x+1)\\
&=x(x-1)(2x-1).
\end{align*}
于是, 当 $x\neq 0$, $x\neq 1$ 时,
\[
C+D(x-1)=2x-1.
\]
令 $x\to 1$ 得到 $C=1$; 令 $x\to 0$ 得到 $1-D=C-D=-1$, $D=2$.

综上所述有 $A=-1$, $B=2$, $C=1$, $D=2$. 积分过程同解法 1.
\end{solution}

\begin{example}\label{ex:5.3.7}
求 $\int\frac{\mathrm{d}x}{x^{4}+1}$.
\end{example}

\begin{solution}[解法 1]
\begin{align*}
\int\frac{\mathrm{d}x}{x^{4}+1}&=\int\frac{\mathrm{d}x}{(x^{2}+1)^{2}-(\sqrt{2}x)^{2}}
=\int\frac{\mathrm{d}x}{(x^{2}+\sqrt{2}x+1)(x^{2}-\sqrt{2}x+1)}\\
&=\int\left(\frac{\frac{\sqrt{2}}{4}x+\frac{1}{2}}{x^{2}+\sqrt{2}x+1}+\frac{-\frac{\sqrt{2}}{4}x+\frac{1}{2}}{x^{2}-\sqrt{2}x+1}\right)\mathrm{d}x\\
&=\frac{1}{4\sqrt{2}}\int\left(\frac{2x+\sqrt{2}}{x^{2}+\sqrt{2}x+1}-\frac{2x-\sqrt{2}}{x^{2}-\sqrt{2}x+1}\right)\mathrm{d}x\\
&\quad+\frac{1}{4}\int\left[\frac{1}{\left(x+\frac{\sqrt{2}}{2}\right)^{2}+\left(\frac{1}{\sqrt{2}}\right)^{2}}
+\frac{1}{\left(x-\frac{\sqrt{2}}{2}\right)^{2}+\left(\frac{1}{\sqrt{2}}\right)^{2}}\right]\mathrm{d}x\\
&=\frac{1}{4\sqrt{2}}\ln\frac{x^{2}+\sqrt{2}x+1}{x^{2}-\sqrt{2}x+1}
+\frac{1}{2\sqrt{2}}\left[\arctan(\sqrt{2}x+1)+\arctan(\sqrt{2}x-1)\right]+C.
\end{align*}
\end{solution}

\begin{solution}[解法 2]
因为
\begin{align*}
\int\frac{x^{2}+1}{x^{4}+1}\mathrm{d}x&=\int\frac{1+\frac{1}{x^{2}}}{x^{2}+\frac{1}{x^{2}}}\mathrm{d}x
=\int\frac{\mathrm{d}\left(x-\frac{1}{x}\right)}{\left(x-\frac{1}{x}\right)^{2}+2}\\
&=\frac{1}{\sqrt{2}}\arctan\frac{x-\frac{1}{x}}{\sqrt{2}}+C_{1},
\end{align*}
\begin{align*}
\int\frac{1-x^{2}}{x^{4}+1}\mathrm{d}x&=\int\frac{\frac{1}{x^{2}}-1}{x^{2}+\frac{1}{x^{2}}}\mathrm{d}x
=-\int\frac{\mathrm{d}\left(x+\frac{1}{x}\right)}{\left(x+\frac{1}{x}\right)^{2}-2}\\
&=-\frac{1}{2\sqrt{2}}\int\left[\frac{1}{\left(x+\frac{1}{x}\right)-\sqrt{2}}
-\frac{1}{\left(x+\frac{1}{x}\right)+\sqrt{2}}\right]\mathrm{d}\left(x+\frac{1}{x}\right)\\
&=\frac{1}{2\sqrt{2}}\ln\frac{x+\frac{1}{x}+\sqrt{2}}{x+\frac{1}{x}-\sqrt{2}}+C_{2}
=\frac{1}{2\sqrt{2}}\ln\frac{x^{2}+\sqrt{2}x+1}{x^{2}-\sqrt{2}x+1}+C_{2},
\end{align*}
所以
\begin{align*}
\int\frac{\mathrm{d}x}{x^{4}+1}&=\frac{1}{2}\int\left(\frac{1-x^{2}}{x^{4}+1}+\frac{x^{2}+1}{x^{4}+1}\right)\mathrm{d}x\\
&=\frac{1}{4\sqrt{2}}\ln\frac{x^{2}+\sqrt{2}x+1}{x^{2}-\sqrt{2}x+1}+\frac{1}{2\sqrt{2}}\arctan\frac{x-\frac{1}{x}}{\sqrt{2}}+C_{3}\\
&=\frac{1}{4\sqrt{2}}\ln\frac{x^{2}+\sqrt{2}x+1}{x^{2}-\sqrt{2}x+1}
+\frac{1}{2\sqrt{2}}\left[\arctan(\sqrt{2}x+1)+\arctan(\sqrt{2}x-1)\right]+C,
\end{align*}
其中, 因为
\begin{align*}
\tan\left[\arctan(\sqrt{2}x+1)+\arctan(\sqrt{2}x-1)\right]
&=\frac{(\sqrt{2}x+1)+(\sqrt{2}x-1)}{1-(2x^{2}-1)}=\frac{\sqrt{2}x}{1-x^{2}},\\
\tan\left(\arctan\frac{x-\frac{1}{x}}{\sqrt{2}}\right)&=\frac{x-\frac{1}{x}}{\sqrt{2}}=\frac{x^{2}-1}{\sqrt{2}x}\\
&=-\cot\left[\arctan(\sqrt{2}x+1)+\arctan(\sqrt{2}x-1)\right],
\end{align*}
所以
\[
\arctan\frac{x-\frac{1}{x}}{\sqrt{2}}=\frac{\pi}{2}+\left[\arctan(\sqrt{2}x+1)+\arctan(\sqrt{2}x-1)\right].
\]
\end{solution}

代数学关于有理函数的部分分式分解定理从理论上证明了任何有理函数的不定积分都为初等函数. 但该有理函数的分母 $Q(x)$, 当它的次数很高时, 具体分解因子是十分困难的. 即使 $Q(x)$ 能因式分解, 用待定系数法将此有理函数化为最简分式之和的计算也会十分烦琐.

现在来介绍几类函数的不定积分, 它们可化为有理函数的不定积分. 因此, 理论上这些函数的不定积分是可以积出来的.

由 $u(x)$, $v(x)$ 及常数经过有限次四则运算所得到的函数称为关于 $u(x)$, $v(x)$ 的有理式, 并表示为 $R(u(x),v(x))=\frac{P(u(x),v(x))}{Q(u(x),v(x))}$, 其中 $P(x,y)=\sum\limits_{i=1}^{n}\sum\limits_{j=1}^{m}a_{ij}x^{i}y^{j}$ ($a_{ij}$ 为实系数, $x,y\in\mathbb{R}$) 为实二元多项式, $Q(x,y)$ 也为实二元多项式.

\subsection*{1. 三角函数有理式的不定积分}

\[
\int R(\sin x,\cos x)\mathrm{d}x
\xlongequal{t=\tan\frac{x}{2}}\int R\left(\frac{2t}{1+t^{2}},\frac{1-t^{2}}{1+t^{2}}\right)\frac{2}{1+t^{2}}\mathrm{d}t.
\]
等式右端是 $t$ 的有理函数的不定积分, 因此, 理论上它可以被积出来, 成为一个初等函数.

正因为如此, 我们称 $t=\tan\frac{x}{2}$ 为万能变换. 其中
\begin{gather*}
x=2\arctan t,\quad \mathrm{d}x=\frac{2}{1+t^{2}}\mathrm{d}t,\\
\sin x=\frac{2\sin\frac{x}{2}\cos\frac{x}{2}}{\sin^{2}\frac{x}{2}+\cos^{2}\frac{x}{2}}
=\frac{2\tan\frac{x}{2}}{\tan^{2}\frac{x}{2}+1}=\frac{2t}{1+t^{2}},\\
\cos x=\frac{\cos^{2}\frac{x}{2}-\sin^{2}\frac{x}{2}}{\sin^{2}\frac{x}{2}+\cos^{2}\frac{x}{2}}
=\frac{1-\tan^{2}\frac{x}{2}}{1+\tan^{2}\frac{x}{2}}=\frac{1-t^{2}}{1+t^{2}}.
\end{gather*}

上面所用的万能变换 $t=\tan\frac{x}{2}$ 对三角函数有理式的不定积分理论上总是有效的. 但如果右边积分中被积函数的分母的次数太高, 往往难以积出来. 此时 ``万能'' 就成为一个虚名了.

另外, 还有几个可化为有理函数的不定积分的变换, 有时这些有理函数分母的幂次较低, 积分比较容易.

(1) 若 $R(-\sin x,\cos x)=-R(\sin x,\cos x)$, 令 $t=\cos x$;

(2) 若 $R(\sin x,-\cos x)=-R(\sin x,\cos x)$, 令 $t=\sin x$;

(3) 若 $R(-\sin x,-\cos x)=R(\sin x,\cos x)$, 令 $t=\tan x$.

\begin{example}\label{ex:5.3.8}
求 $\int\frac{1+\sin x}{\sin x(1+\cos x)}\mathrm{d}x$.
\end{example}

\begin{solution}
\begin{align*}
&\int\frac{1+\sin x}{\sin x(1+\cos x)}\mathrm{d}x\\
&\xlongequal{t=\tan\frac{x}{2}}\int\frac{1+\frac{2t}{1+t^{2}}}{\frac{2t}{1+t^{2}}\left(1+\frac{1-t^{2}}{1+t^{2}}\right)}\cdot\frac{2}{1+t^{2}}\mathrm{d}t\\
&=\int\frac{1+t^{2}+2t}{t(1+t^{2}+1-t^{2})}\mathrm{d}t=\frac{1}{2}\int\left(t+2+\frac{1}{t}\right)\mathrm{d}t\\
&=\frac{1}{2}\left(\frac{t^{2}}{2}+2t+\ln|t|\right)+C\\
&=\frac{1}{4}\tan^{2}\frac{x}{2}+\tan\frac{x}{2}+\frac{1}{2}\ln\left|\tan\frac{x}{2}\right|+C.
\end{align*}
\end{solution}

\begin{example}\label{ex:5.3.9}
求 $\int\frac{\mathrm{d}x}{a^{2}\sin^{2}x+b^{2}\cos^{2}x}$ ($a>0$, $b>0$).
\end{example}

\begin{solution}
\begin{align*}
\int\frac{\mathrm{d}x}{a^{2}\sin^{2}x+b^{2}\cos^{2}x}
&=\int\frac{\frac{1}{\cos^{2}x}}{a^{2}\frac{\sin^{2}x}{\cos^{2}x}+b^{2}}\mathrm{d}x\\
&=\int\frac{\mathrm{d}\tan x}{a^{2}\tan^{2}x+b^{2}}
=\frac{1}{ab}\arctan\left(\frac{a}{b}\arctan x\right)+C.
\end{align*}
\end{solution}

\begin{example}\label{ex:5.3.10}
求 $\int\frac{\mathrm{d}x}{\sin^{2}x\cos x}$.
\end{example}

\begin{solution}
\begin{align*}
\int\frac{\mathrm{d}x}{\sin^{2}x\cos x}&=\int\frac{\mathrm{d}\sin x}{\sin^{2}x(1-\sin^{2}x)}
\xlongequal{t=\sin x}\int\frac{\mathrm{d}t}{t^{2}(1-t^{2})}\\
&=\int\left(\frac{1}{t^{2}}+\frac{1}{1-t^{2}}\right)\mathrm{d}t=-\frac{1}{t}+\frac{1}{2}\ln\frac{1+t}{1-t}+C\\
&=-\frac{1}{\sin x}+\frac{1}{2}\ln\frac{1+\sin x}{1-\sin x}+C.
\end{align*}
\end{solution}

\begin{example}\label{ex:5.3.11}
求 $\int\frac{\sin^{4}x}{\cos^{2}x}\mathrm{d}x$.
\end{example}

\begin{solution}[解法 1]
\begin{align*}
\int\frac{\sin^{4}x}{\cos^{2}x}\mathrm{d}x&=\int\frac{\sin^{4}x}{(\sin^{2}x+\cos^{2}x)^{2}}\mathrm{d}\tan x
=\int\frac{\tan^{4}x}{(\tan^{2}x+1)^{2}}\mathrm{d}\tan x\\
&\xlongequal{t=\tan x}\int\frac{t^{4}}{(t^{2}+1)^{2}}\mathrm{d}t
=\int\frac{(t^{4}+2t^{2}+1)-2(t^{2}+1)+1}{(t^{2}+1)^{2}}\mathrm{d}t\\
&=\int\left[1-\frac{1}{1+t^{2}}+\frac{1}{(1+t^{2})^{2}}\right]\mathrm{d}t\\
&\xlongequal{\text{递推公式}}t-2\arctan t+\left(\frac{1}{2}\cdot\frac{t}{1+t^{2}}+\frac{1}{2}\arctan t\right)+C\\
&=t-\frac{3}{2}\arctan t+\frac{t}{2(1+t^{2})}+C\\
&=\tan x-\frac{3}{2}x+\frac{1}{2}\cdot\frac{\tan x}{\sec^{2}x}+C\\
&=\tan x-\frac{3}{2}x+\frac{1}{4}\sin 2x+C.
\end{align*}
\end{solution}

\begin{solution}[解法 2]
\begin{align*}
\int\frac{\sin^{4}x}{\cos^{2}x}\mathrm{d}x&\xlongequal{t=\tan x}\int\frac{t^{4}}{(t^{2}+1)^{2}}\mathrm{d}t\\
&=\int\left[1-\frac{1}{t^{2}+1}-\frac{t^{2}}{(t^{2}+1)^{2}}\right]\mathrm{d}t\\
&=t-\arctan t+\frac{1}{2}\int t\,\mathrm{d}\frac{1}{t^{2}+1}\\
&=t-\arctan t+\frac{t}{2(t^{2}+1)}-\frac{1}{2}\int\frac{\mathrm{d}t}{t^{2}+1}\\
&=t-\frac{3}{2}\arctan t+\frac{t}{2(t^{2}+1)}+C\\
&=\tan x-\frac{3}{2}x+\frac{1}{2}\cdot\frac{\tan x}{\sec^{2}x}+C\\
&=\tan x-\frac{3}{2}x+\frac{1}{4}\sin 2x+C.
\end{align*}
\end{solution}

\begin{example}\label{ex:5.3.12}
求 Poisson 积分 $\int\frac{1-r^{2}}{1-2r\cos x+r^{2}}\mathrm{d}x$, $0<r<1$, $|x|<\pi$.
\end{example}

\begin{solution}
作万能变换 $t=\tan\frac{x}{2}$, 得到
\begin{align*}
&\int\frac{1-r^{2}}{1-2r\cos x+r^{2}}\mathrm{d}x\\
&=\int\frac{1-r^{2}}{1-2r\frac{1-t^{2}}{1+t^{2}}+r^{2}}\cdot\frac{2}{1+t^{2}}\mathrm{d}t\\
&=2(1-r^{2})\int\frac{\mathrm{d}t}{1+t^{2}-2r+2rt^{2}+r^{2}+r^{2}t^{2}}\\
&=2(1-r^{2})\int\frac{\mathrm{d}t}{(1-r)^{2}+(1+r)^{2}t^{2}}\\
&=\frac{2(1-r^{2})}{(1+r)^{2}}\int\frac{\mathrm{d}t}{\left(\frac{1-r}{1+r}\right)^{2}+t^{2}}\\
&=\frac{2(1-r^{2})}{(1+r)^{2}}\cdot\frac{1}{\frac{1-r}{1+r}}\arctan\frac{t}{\frac{1-r}{1+r}}+C\\
&=2\arctan\left(\frac{1+r}{1-r}\tan\frac{x}{2}\right)+C.
\end{align*}
\end{solution}

\subsection*{2. 某些带无理根式的不定积分}

设 $R(x,y_{1},\cdots,y_{k})$ 为 $k+1$ 元实有理函数.

(1) $\int R\left(x,\sqrt[n]{\frac{ax+b}{cx+d}},\cdots,\sqrt[m]{\frac{ax+b}{cx+d}}\right)\mathrm{d}x$ 型不定积分, 其中 $a,b,c,d\in\mathbb{R}$, $ad-bc\neq 0$, $n,\cdots,m\in\mathbb{N}$. 令
\[
t=\left(\frac{ax+b}{cx+d}\right)^{\frac{1}{s}},
\]
这里 $s$ 为 $n,\cdots,m$ 的最小公倍数, 则
\begin{gather*}
\frac{ax+b}{cx+d}=t^{s},\\
ax+b=ct^{s}x+t^{s}d,\\
x=\frac{dt^{s}-b}{a-ct^{s}},\\
\mathrm{d}x=\frac{st^{s-1}(ad-bc)}{(a-ct^{s})^{2}}\mathrm{d}t.
\end{gather*}

(2) $\int R(x,\sqrt{ax^{2}+bx+c})\mathrm{d}x$ 型不定积分, 其中当 $a>0$ 时, $b^{2}-4ac\neq 0$; 当 $a<0$ 时, $b^{2}-4ac>0$. 由于
\[
ax^{2}+bx+c=a\left[\left(x+\frac{b}{2a}\right)^{2}+\frac{4ac-b^{2}}{4a^{2}}\right],
\]
若记 $u=x+\frac{b}{2a}$, $k^{2}=\left|\frac{4ac-b^{2}}{4a^{2}}\right|$, 则上述二次三项式必属于以下三种情形之一:
\[
|a|(u^{2}+k^{2}),\quad |a|(u^{2}-k^{2}),\quad |a|(k^{2}-u^{2}).
\]
因此, 上面的不定积分转化为下面三种类型不定积分之一:
\[
\int R(u,\sqrt{u^{2}\pm k^{2}})\mathrm{d}u,\quad \int R(u,\sqrt{k^{2}-u^{2}})\mathrm{d}u.
\]
当分别令 $u=k\tan t$, $u=k\sec t$, $u=k\sin t$ 后, 它们都可化为三角有理式的不定积分. 特别地, 还有下面几种情况.

① $a>0$, 作变换 $\sqrt{ax^{2}+bx+c}=t-\sqrt{a}x$ (或 $\sqrt{a}x\pm t$), 于是,
\begin{gather*}
t=\sqrt{ax^{2}+bx+c}+\sqrt{a}x,\\
ax^{2}+bx+c=t^{2}-2\sqrt{a}tx+ax^{2},
\end{gather*}
两边消去 $ax^{2}$, 移项得
\begin{gather*}
(b+2\sqrt{a}t)x=t^{2}-c,\\
x=\frac{t^{2}-c}{b+2\sqrt{a}t},\\
\mathrm{d}x=\frac{2t(b+2\sqrt{a}t)-2\sqrt{a}(t^{2}-c)}{(b+2\sqrt{a}t)^{2}}\mathrm{d}t\\
=\frac{2bt+2\sqrt{a}t^{2}+2\sqrt{ac}}{(b+2\sqrt{a}t)^{2}}\mathrm{d}t.
\end{gather*}

② $c>0$, 作变换 $\sqrt{ax^{2}+bx+c}=tx+\sqrt{c}$ (或 $tx-\sqrt{c}$), 则
\begin{gather*}
t=\frac{\sqrt{ax^{2}+bx+c}-\sqrt{c}}{x},\\
ax^{2}+bx+c=t^{2}x^{2}+2\sqrt{c}tx+c.
\end{gather*}
两边消去 $c$, 约去 $x$ 得
\begin{gather*}
ax+b=t^{2}x+2\sqrt{c}t,\\
x=\frac{2\sqrt{c}t-b}{a-t^{2}},\\
\mathrm{d}x=\frac{2a\sqrt{c}-2bt+2\sqrt{c}t^{2}}{(a-t^{2})^{2}}\mathrm{d}t.
\end{gather*}

③ $b^{2}-4ac>0$, 则 $ax^{2}+bx+c=0$ 有两个相异的实根 $\lambda,\mu$. 于是
\[
ax^{2}+bx+c=a(x-\lambda)(x-\mu).
\]
作变换
\[
\sqrt{ax^{2}+bx+c}=t(x-\lambda),
\]
两边平方得
\begin{gather*}
a(x-\lambda)(x-\mu)=ax^{2}+bx+c=t^{2}(x-\lambda)^{2},\\
a(x-\mu)=t^{2}(x-\lambda),\\
x=\frac{a\mu-\lambda t^{2}}{a-t^{2}},\\
\mathrm{d}x=\frac{2a(\mu-\lambda)t}{(a-t^{2})^{2}}\mathrm{d}t.
\end{gather*}

④ $b^{2}-4ac=0$, 当 $a>0$, 则 $\sqrt{ax^{2}+bx+c}=\sqrt{a(x-\lambda)^{2}}=\sqrt{a}(x-\lambda)$ 或 $\sqrt{a}(\lambda-x)$, 此时, $R(x,\sqrt{ax^{2}+bx+c})$ 为 $x$ 的有理式;

当 $a<0$ 时, $\sqrt{ax^{2}+bx+c}=\sqrt{a(x-\lambda)^{2}}$ 对 $x\neq\lambda$ 无意义.

⑤ $b^{2}-4ac<0$, 则 $ax^{2}+bx+c=a\left[\left(x+\frac{b}{2a}\right)^{2}+\frac{4ac-b^{2}}{4a^{2}}\right]$ 与 $a$ 同号.

当 $a>0$ 时, 由 ① 作变换 $x=\frac{t^{2}-c}{b+2\sqrt{a}t}$;

当 $a<0$ 时, $\sqrt{ax^{2}+bx+c}$ 无意义.

\begin{example}\label{ex:5.3.13}
求 $\int\frac{\mathrm{d}x}{x\sqrt{x^{2}-2x-3}}$.
\end{example}

\begin{solution}[解法 1]
\begin{align*}
\int\frac{\mathrm{d}x}{x\sqrt{x^{2}-2x-3}}&=\int\frac{\mathrm{d}(x-1)}{x\sqrt{(x-1)^{2}-4}}\\
&\xlongequal{u=x-1}\int\frac{\mathrm{d}u}{(u+1)\sqrt{u^{2}-4}}
\xlongequal{u=2\sec\theta}\int\frac{2\sec\theta\cdot\tan\theta}{(2\sec\theta+1)\cdot 2\tan\theta}\mathrm{d}\theta\\
&=\int\frac{\mathrm{d}\theta}{2+\cos\theta}
\xlongequal{t=\tan\frac{\theta}{2}}\int\frac{\frac{2}{1+t^{2}}}{2+\frac{1-t^{2}}{1+t^{2}}}\mathrm{d}t
=\int\frac{2}{t^{2}+3}\mathrm{d}t=\frac{2}{\sqrt{3}}\arctan\frac{t}{\sqrt{3}}+C\\
&=\frac{2}{\sqrt{3}}\arctan\left(\frac{1}{\sqrt{3}}\tan\frac{\theta}{2}\right)+C
=\frac{2}{\sqrt{3}}\arctan\frac{\sqrt{x^{2}-2x-3}}{\sqrt{3}(x+1)}+C.
\end{align*}
其中
\begin{align*}
\tan\frac{\theta}{2}&=\frac{2\sin\frac{\theta}{2}\cos\frac{\theta}{2}}{2\cos^{2}\frac{\theta}{2}}
=\frac{\sin\theta}{1+\cos\theta}=\frac{\tan\theta}{\sec\theta+1}\\
&=\frac{\sqrt{\left(\frac{u}{2}\right)^{2}-1}}{\frac{u}{2}+1}
=\frac{\sqrt{\left(\frac{x-1}{2}\right)^{2}-1}}{\frac{x-1}{2}+1}
=\frac{\sqrt{x^{2}-2x-3}}{x+1}.
\end{align*}
\end{solution}

\begin{solution}[解法 2]
$a=1>0$, 令 $\sqrt{x^{2}-2x-3}=x-t$, 两边平方得到
\begin{gather*}
x^{2}-2x-3=x^{2}-2tx+t^{2},\\
-2x-3=-2tx+t^{2},\\
x=\frac{t^{2}+3}{2(t-1)},\\
\mathrm{d}x=\frac{2t(t-1)-(t^{2}+3)}{2(t-1)^{2}}\mathrm{d}t=\frac{t^{2}-2t-3}{2(t-1)^{2}}\mathrm{d}t,\\
\sqrt{x^{2}-2x-3}=x-t=\frac{t^{2}+3}{2(t-1)}-t=\frac{-(t^{2}-2t-3)}{2(t-1)}.
\end{gather*}
于是,
\begin{align*}
\int\frac{\mathrm{d}x}{x\sqrt{x^{2}-2x-3}}
&=\int\frac{\frac{t^{2}-2t-3}{2(t-1)^{2}}}{\frac{t^{2}+3}{2(t-1)}\cdot\frac{-(t^{2}-2t-3)}{2(t-1)}}\mathrm{d}t\\
&=-2\int\frac{\mathrm{d}t}{t^{2}+3}=-\frac{2}{\sqrt{3}}\arctan\frac{t}{\sqrt{3}}+C_{1}\\
&=\frac{2}{\sqrt{3}}\arctan\frac{\sqrt{x^{2}-2x-3}-x}{\sqrt{3}}+C_{1}\\
&=\frac{2}{\sqrt{3}}\left(\arctan\frac{\sqrt{x^{2}-2x-3}}{\sqrt{3}(x+1)}-\frac{\pi}{3}\right)+C_{1}\\
&=\frac{2}{\sqrt{3}}\arctan\frac{\sqrt{x^{2}-2x-3}}{\sqrt{3}(x+1)}+C.
\end{align*}
\end{solution}

读者必须知道, 存在初等函数, 它的不定积分 ``积不出来'', 即原函数不能用初等函数表达. 例如:
\[
\int\mathrm{e}^{\pm x^{2}}\mathrm{d}x,\quad \int\sin x^{2}\,\mathrm{d}x,\quad \int\cos x^{2}\,\mathrm{d}x,\quad
\mathrm{six}=\int\frac{\sin x}{x}\mathrm{d}x\ (\text{积分正弦}),
\]
\[
\mathrm{cix}=\int\frac{\cos x}{x}\mathrm{d}x\ (\text{积分余弦}),\quad
\mathrm{Lix}=\int\frac{\mathrm{d}x}{\ln x}\ (\text{积分对数}),
\]
Legendre 形式下的第一、二、三类椭圆积分 ($0<k<1$)
\[
\int\frac{\mathrm{d}x}{\sqrt{1-k^{2}\sin^{2}x}},\quad \int\sqrt{1-k^{2}\sin^{2}x}\,\mathrm{d}x,\quad
\int\frac{\mathrm{d}x}{(1+h\sin^{2}x)\sqrt{1-k^{2}\sin^{2}x}}.
\]
虽然它们都存在 (参阅第 6 章微积分基本定理~\ref{thm:6.3.2}), 但不能用初等函数表达 (证明非常艰难, Liouville 于 1835 年作出证明). 因此, 初等函数的原函数不一定是初等函数.

\begin{exblock}{练习题 5.3}
  \item 求下列有理函数的不定积分:
  \begin{enumerate}[label=(\arabic*)]
    \item $\int\frac{x^{3}}{x-1}\mathrm{d}x$;
    \item $\int\frac{x-2}{x^{2}-7x+12}\mathrm{d}x$;
    \item $\int\frac{\mathrm{d}x}{1+x^{3}}$;
    \item $\int\frac{\mathrm{d}x}{1+x^{4}}$;
    \item $\int\frac{\mathrm{d}x}{(x-1)(x^{2}+1)^{2}}$;
    \item $\int\frac{x-2}{(2x^{2}+2x+1)^{2}}\mathrm{d}x$;
    \item $\int\frac{5x+6}{x^{2}+x+1}\mathrm{d}x$;
    \item $\int\frac{x}{2x^{2}-3x-2}\mathrm{d}x$;
    \item $\int\frac{2x^{2}+1}{(x+3)(x-1)(x-4)}\mathrm{d}x$;
    \item $\int\frac{\mathrm{d}x}{2x^{3}+3x^{2}+x}$;
    \item $\int\frac{x^{2}}{x^{3}+5x^{2}+8x+4}\mathrm{d}x$;
    \item $\int\frac{4x+3}{(x-2)^{3}}\mathrm{d}x$;
    \item $\int\frac{x^{2}}{(x+2)^{2}(x+4)^{2}}\mathrm{d}x$;
    \item $\int\frac{8x^{3}+7}{(x+1)(2x+1)^{3}}\mathrm{d}x$;
    \item $\int\frac{\mathrm{d}x}{(x+1)(x+2)^{2}(x+3)^{3}}$;
    \item $\int\frac{\mathrm{d}x}{(x^{2}+9)^{3}}$;
    \item $\int\frac{1+x^{2}}{1+x^{4}}\mathrm{d}x$;
    \item $\int\frac{\mathrm{d}x}{x^{4}(1+x^{2})}$;
    \item $\int\frac{x-5}{x^{3}-3x^{2}+4}\mathrm{d}x$;
    \item $\int\frac{x^{7}}{x^{4}+2}\mathrm{d}x$.
  \end{enumerate}

  \item 求下列不定积分:
  \begin{enumerate}[label=(\arabic*)]
    \item $\int\frac{\mathrm{d}x}{\sqrt{x^{2}+x}}$;
    \item $\int\frac{1}{x^{2}}\sqrt{\frac{1-x}{1+x}}\mathrm{d}x$;
    \item $\int\frac{\mathrm{d}x}{\sqrt{x}(1+\sqrt[3]{x})}$;
    \item $\int\sqrt{\frac{x-a}{b-x}}\mathrm{d}x$, $a<x<b$;
    \item $\int\frac{\mathrm{d}x}{1+\sqrt{x}}$;
    \item $\int\frac{\mathrm{d}x}{x\sqrt{x^{2}-1}}$;
    \item $\int\frac{\mathrm{d}x}{1+\sqrt[3]{1+x}}$;
    \item $\int\frac{\sqrt[3]{x}}{x(\sqrt{x}+\sqrt[3]{x})}\mathrm{d}x$;
    \item $\int\frac{x}{\sqrt{x+1}+\sqrt[3]{x+1}}\mathrm{d}x$;
    \item $\int\frac{\sqrt{1+x^{2}}}{2+x^{2}}\mathrm{d}x$;
    \item $\int\frac{\mathrm{d}x}{x\sqrt{1+x^{2}}}$;
    \item $\int\frac{x^{3}}{\sqrt{1+2x^{2}}}\mathrm{d}x$;
    \item $\int\sqrt{\frac{1-\sqrt{x}}{1+\sqrt{x}}}\mathrm{d}x$;
    \item $\int\frac{\mathrm{d}x}{x\sqrt{a^{2}-x^{2}}}$, $|x|<a$, $x\neq 0$;
    \item $\int\frac{x^{2}}{\sqrt{1+x-x^{2}}}\mathrm{d}x$;
    \item $\int\frac{\mathrm{d}x}{(x+a)^{2}(x+b)^{3}}$.
  \end{enumerate}

  \item 求下列不定积分:
  \begin{enumerate}[label=(\arabic*)]
    \item $\int\frac{\mathrm{d}x}{5-3\cos x}$;
    \item $\int\frac{\mathrm{d}x}{2+\sin^{2}x}$;
    \item $\int\frac{\mathrm{d}x}{1+\tan x}$;
    \item $\int\frac{\tan x}{1+\tan x+\tan^{2}x}\mathrm{d}x$;
    \item $\int\frac{\sin^{3}x}{\cos^{4}x}\mathrm{d}x$;
    \item $\int\frac{\mathrm{d}x}{\cos^{4}x\sin x}$;
    \item $\int\frac{1-\tan x}{1+\tan x}\mathrm{d}x$. (提示: $\frac{1-\tan x}{1+\tan x}=\frac{\cos x-\sin x}{\cos x+\sin x}$ 或 $u=\tan x$.)
  \end{enumerate}
\end{exblock}
